\documentclass[11pt,letterpaper]{article}

\usepackage[margin=1.1in]{geometry}
\usepackage{amsmath,amssymb,amsthm}
\usepackage{booktabs}
\usepackage{array}
\usepackage{microtype}
\usepackage{hyperref}
\usepackage{xcolor}
\usepackage{algorithm}
\usepackage{algpseudocode}
\usepackage{caption}

\hypersetup{colorlinks=true, linkcolor=blue!60!black, citecolor=blue!60!black, urlcolor=blue!60!black}

\newtheorem{theorem}{Theorem}
\newtheorem{proposition}[theorem]{Proposition}
\newtheorem{definition}[theorem]{Definition}
\newtheorem{remark}[theorem]{Remark}

\title{\textbf{Exact Equilibrium Solutions for the Board Game \emph{Football Strategy}:
A Full-Game Backward-Induction Solver}\\[6pt]
\large An Extension of McCulloch's Game-Theoretic Agent}

\author{George Kyrollos}

\date{}

\begin{document}

\maketitle

\begin{abstract}
McCulloch~\cite{mcculloch} demonstrated that each individual scrimmage play in Avalon
Hill's board game \emph{Football Strategy} constitutes a two-player zero-sum matrix game,
and designed an agent that solves each play in isolation using mixed-strategy Nash
equilibrium.  To make the agent sensitive to game context---field position, down, distance,
score, and clock---McCulloch augmented the raw yardage payoffs with hand-crafted heuristic
bonuses.  We replace that heuristic layer entirely.  By modeling the complete game as a
finite-horizon two-player zero-sum stochastic game and applying exact backward induction,
every matrix-game entry becomes the true equilibrium continuation value of the resulting
game state.  The solver determines optimal play across all ${\sim}30$ million reachable
states of the fourth quarter, selecting between punt, field goal, and scrimmage plays based
on actual win probability rather than proxy rewards.  Each team's three timeouts per half are modeled as a pure-strategy sequential
sub-game resolved after each play, so the solution reflects optimal timeout usage.  We describe the
mathematical formulation, state-space representation, solution algorithm, and computational
architecture, and compare the approach explicitly to McCulloch's method.
\end{abstract}

%────────────────────────────────────────────────────────────────────────────────
\section{Introduction}
\label{sec:intro}
%────────────────────────────────────────────────────────────────────────────────

\emph{Football Strategy} (Avalon Hill, 1969) is a two-player board game that abstracts
American football into a sequence of simultaneous-move interactions.  On each scrimmage
play the offense secretly selects one of up to 22 actions while the defense simultaneously
selects one of ten cards (A--J); the intersection on a printed chart determines the
result---yardage, an incomplete pass, a fumble, an interception, a penalty, or a score.
Because both players choose without observing the opponent's action, the play is exactly a
$|O|\times 10$ two-player zero-sum matrix game.

McCulloch~\cite{mcculloch} was the first to exploit this structure algorithmically.  His
agent solves each matrix game using linear programming to obtain a mixed-strategy Nash
equilibrium.  The key challenge is assigning \emph{utility} to each chart outcome: a raw
yardage gain of $+7$ is not inherently better or worse than a gain of $+4$ without knowing
the down, distance, score, clock, and field position.  McCulloch addressed this with a set
of hand-tuned bonuses---for first downs, touchdowns, turnovers, and fourth-down
conversions---that effectively map game states to a scalar reward.  While principled in
spirit, these heuristics are necessarily approximate and do not guarantee that the agent's
play is optimal for the actual game of \emph{Football Strategy}.

Our contribution is to eliminate the heuristic layer completely.  Instead of assigning an
ad-hoc scalar to each outcome, we assign it the \emph{exact equilibrium continuation value}
of the full game state it produces.  This requires solving ${\sim}30$ million linked matrix
games simultaneously via backward induction over the complete finite game tree.  The result
is a provably optimal strategy---in the sense of the minimax theorem for finite two-player
zero-sum stochastic games---for the fourth quarter of \emph{Football Strategy}.

%────────────────────────────────────────────────────────────────────────────────
\section{McCulloch's Approach and Its Limitations}
\label{sec:mcculloch}
%────────────────────────────────────────────────────────────────────────────────

McCulloch's agent operates as follows.  At each scrimmage state it constructs a payoff
matrix $A \in \mathbb{R}^{|O|\times 10}$ where entry $A_{a,b}$ is the utility assigned to
the outcome produced by offensive action $a$ and defensive card $b$.  The matrix game
$\max_{\sigma_O}\min_{\sigma_D}\, \sigma_O^\top A\, \sigma_D$ is then solved via LP to
obtain mixed strategies.

The utility function $A_{a,b}$ is where the approximation lives.  McCulloch converts each
chart result into a scalar by applying additive bonuses:

\begin{itemize}
  \item A first down earns a bonus proportional to the fraction of the field gained.
  \item A touchdown earns a large positive bonus scaled by the score context.
  \item An interception or fumble lost earns a large negative bonus.
  \item A fourth-down conversion earns a bonus; failing on fourth down earns a penalty.
  \item Clock and score interact only through the bonus magnitudes, which are set by hand.
\end{itemize}

The limitations of this approach are fundamental rather than incidental:

\paragraph{Heuristics cannot capture non-monotone value.}
The value of gaining 10 yards on first-and-10 at midfield in the final minute while
trailing by 3 is qualitatively different from the value of the same play with ten minutes
remaining.  No fixed additive bonus captures this interaction without effectively
re-solving the game.

\paragraph{Optimal fourth-down decisions require global knowledge.}
Whether to punt, attempt a field goal, or go for it on fourth down depends on score,
clock, and field position in ways that require comparing continuation values under each
option.  McCulloch's agent applies fixed cutoffs; the correct decision is state-dependent
and only determinable via backward induction.

\paragraph{End-of-half clock management is inaccessible.}
The strategic value of running out the clock, spiking the ball, or calling a timeout is
zero under heuristics but central to real play.  Capturing it requires the value function
to be defined over clock states.

%────────────────────────────────────────────────────────────────────────────────
\section{Formal Game Model}
\label{sec:model}
%────────────────────────────────────────────────────────────────────────────────

\subsection{Terminal Utility}

We model \emph{Football Strategy} as a finite-horizon two-player zero-sum stochastic game.
Let $\delta \in \mathbb{Z}$ denote the score differential (Team~1 minus Team~2) at the end
of the game.  The terminal utility is
\[
  U(\delta) = \operatorname{sgn}(\delta) \in \{-1, 0, +1\},
\]
so the objective is to maximize win probability, not expected point margin.  This is the
correct objective for a game whose outcome is win, loss, or draw.

\subsection{State Representation}

A complete game state must encode every quantity that affects future optimal play.  We
use a \emph{symmetric possession} convention: the state is always written from the
perspective of the team currently in possession of the ball, so field position $x \in
[1,99]$ is always measured in yards from the current offense's own goal line
($x = 0$: own goal; $x = 100$: opponent's goal).

\begin{definition}[Scrimmage state]
A scrimmage state is a tuple
\[
  s = (q,\, \tau,\, x,\, d,\, y,\, \delta,\, h,\, \theta_1,\, \theta_2),
\]
where:
\begin{itemize}
  \item $q \in \{1,2,3,4\}$ is the quarter;
  \item $\tau \in \{0,1,\ldots,60\}$ is the number of 15-second ticks remaining in the quarter;
  \item $x \in \{1,\ldots,99\}$ is ball position (yards from current offense's own goal);
  \item $d \in \{1,2,3,4\}$ is the current down;
  \item $y \in \{1,\ldots,30\}$ is yards to first down (capped at 30);
  \item $\delta \in \{-35,\ldots,35\}$ is the score differential from the current offense's
        perspective (clipped at $\pm 35$);
  \item $h \in \{-1,0,+1\}$ encodes which team receives the second-half kickoff ($h=0$ in
        Q3/Q4 when this is already resolved); and
  \item $\theta_1,\theta_2 \in \{0,1,2,3\}$ are the remaining timeouts for the current
        offense and defense, respectively.
\end{itemize}
\end{definition}

\noindent The full game starts at $\theta_1 = \theta_2 = 3$.

\noindent In addition to scrimmage states, the game includes \emph{kickoff states}
$(q,\tau,\delta,h)$ and \emph{safety-kick states} $(q,\tau,\delta,h)$.

\paragraph{Possession symmetry.}
Because the state is always expressed from the current offense's viewpoint, a possession
change transforms the state as
\[
  x \;\mapsto\; 100 - x, \qquad
  \delta \;\mapsto\; -\delta, \qquad
  h \;\mapsto\; -h,
\]
and the continuation value is negated (since the current maximizer becomes the minimizer).
This halves the effective state space compared to a possession-explicit encoding.

\subsection{Transition Kernel}

The transition kernel $P(s' \mid s, a, b)$ is determined by:
\begin{enumerate}
  \item \textbf{Chart lookup:} the cell at row $b$ (defense card), column $a$ (offense
        play) of the selected offense chart.
  \item \textbf{Outcome parsing:} chart cells encode one of: a deterministic yardage gain
        or loss (possibly out-of-bounds), an incomplete pass, a fumble, an interception with
        return yards, a long gain (LG), a penalty, a loss-of-down, or an OR-choice between
        two outcomes resolved by the favored team.
  \item \textbf{Game-rule application:} yardage is added to $x$; if $x \geq 100$ a
        touchdown results; if $x \leq 0$ a safety results; first downs reset $d$ and $y$;
        fourth-down failures transfer possession.
  \item \textbf{Clock advancement:} $\tau$ is decremented by the tick cost of the outcome
        type (Table~\ref{tab:ticks}), subject to the two-minute automatic stoppage at
        $\tau = 8$ in Q2 and Q4.
\end{enumerate}

\begin{table}[ht]
\centering
\caption{Time consumed per outcome category (15~sec per tick).}
\label{tab:ticks}
\begin{tabular}{lcc}
\toprule
Outcome type & Time (sec) & Ticks \\
\midrule
In-bounds gain of $0$--$19$ yards & 30 & 2 \\
In-bounds gain of $\geq 20$ yards & 45 & 3 \\
Any loss & 30 & 2 \\
Out-of-bounds play & 15 & 1 \\
Incomplete pass & 15 & 1 \\
Interception & 30 & 2 \\
Fumble & 15 & 1 \\
Penalty & 15 & 1 \\
Kickoff / field goal / punt & 15 & 1 \\
Extra point (PAT) & 0 & 0 \\
\bottomrule
\end{tabular}
\end{table}

\paragraph{Long Gain.}
Chart cells marked \texttt{LG} trigger a probabilistic outcome drawn from the distribution
in Table~\ref{tab:lg}, which was derived from the printed Long Gain table.  LG outcomes
always consume 3 ticks.

\begin{table}[ht]
\centering
\caption{Long Gain distribution.}
\label{tab:lg}
\begin{tabular}{cc}
\toprule
Gain (yards) & Probability \\
\midrule
$+30,\;+35,\;+40,\;+45,\;+50$ & $\tfrac{1}{6}$ each \\
$+60,\;+70,\;\ldots,\;+110$ & $\tfrac{1}{36}$ each \\
\bottomrule
\end{tabular}
\end{table}

\paragraph{Field goals.}
A field goal attempt is legal on any down when $x \geq 55$ (within 45 yards of the
opponent's goal line).  Success probabilities are distance-based (Table~\ref{tab:fg}).

\begin{table}[ht]
\centering
\caption{Field goal success probabilities by distance.}
\label{tab:fg}
\begin{tabular}{cc}
\toprule
Distance to goal line (yards) & $P(\text{good})$ \\
\midrule
$1$--$12$   & $\tfrac{32}{36}$ \\
$13$--$22$  & $\tfrac{28}{36}$ \\
$23$--$32$  & $\tfrac{18}{36}$ \\
$33$--$38$  & $\tfrac{12}{36}$ \\
$39$--$45$  & $\tfrac{2}{36}$  \\
\bottomrule
\end{tabular}
\end{table}

\paragraph{Extra points.}
After every touchdown the scoring team attempts a PAT with $P(\text{good}) = \tfrac{34}{36}$.
PATs consume zero clock time.

%────────────────────────────────────────────────────────────────────────────────
\section{Solution Algorithm}
\label{sec:algorithm}
%────────────────────────────────────────────────────────────────────────────────

\subsection{The Core Idea}

The fundamental observation is that \emph{Football Strategy} is finite: the clock
always ticks forward, so the game must end.  This means we can work backwards from
the final whistle.  At the last possible moment (clock at zero), the outcome is just
win, loss, or draw --- the terminal utility $U(\delta) = \operatorname{sgn}(\delta)$.
One tick earlier, each team chooses optimally given that they know exactly what will
happen at clock zero.  One tick before that, they choose knowing what will happen one
tick from now.  Repeat all the way back to kickoff.  This is backward induction, and
it gives the exact equilibrium of the game.

\subsection{Solving Each Game State}

At each scrimmage state $s$, both teams choose simultaneously and without seeing each
other's decision --- the offense picks a play number, the defense picks a card, and the
result is read from the chart.  This is exactly a finite two-player zero-sum matrix game.
The payoff matrix $A^{(s)}$ has one row per legal offensive action and ten columns (one
per defense card~A--J), with entry
\[
  A^{(s)}_{a,b} = \sum_{s'} P(s'\mid s,a,b)\, V(s'),
\]
the expected game value of the position that results from the play.  Since we solve
from low clock to high clock, $V(s')$ is always already known.

The equilibrium value of state $s$ is the solution to the minimax problem:
\begin{equation}
  \label{eq:minimax}
  V(s) = \max_{\sigma_O}\; \min_{\sigma_D}\;
  \sigma_O^\top A^{(s)} \sigma_D.
\end{equation}
By the minimax theorem, there is always a well-defined value and optimal (possibly
mixed) strategies for both sides.

\subsection{Finding the Equilibrium: Saddle Points and Linear Programming}

For many game states there is a single dominant offensive play --- one that is at
least as good as every other play regardless of what the defense calls.  We call this
a \emph{saddle point}: the maximum of the row-minimum values equals the minimum of
the column-maximum values,
\[
  \max_a \min_b A_{a,b} = \min_b \max_a A_{a,b}.
\]
When this holds, the optimal strategy is pure (no randomisation needed), and we read
the answer directly off the matrix in $O(|O|\cdot|D|)$ time without any LP call.
This is not an approximation --- it is the exact Nash equilibrium, guaranteed by the
saddle-point theorem.  In practice this shortcut fires on roughly 40--60\% of game
states, substantially reducing total solve time.

When no saddle point exists, the equilibrium is a mixed strategy found by solving the
row player's (offense's) LP:
\begin{align*}
  \text{maximize} \quad & v \\
  \text{subject to} \quad
  & \sum_{a} p_a\, A^{(s)}_{a,b} \geq v \quad \forall b, \\
  & \textstyle\sum_{a} p_a = 1, \quad p_a \geq 0.
\end{align*}
The optimal $v^* = V(s)$ and $p^*$ is the offense's equilibrium mixed strategy.
Both teams' strategies are computed using the HiGHS solver via
\texttt{scipy.optimize.linprog}.

\subsection{Backward Induction and Forward Reachability}

Because the clock only moves forward, all states at time $\tau$ depend only on
states at $\tau' < \tau$ --- there are no circular dependencies.  This means we can
solve the whole game in a single sweep: process the states at $\tau=1$ first (only
one tick of game remaining), then $\tau=2$, all the way up to $\tau=60$.  Within
each quarter, every state at a given $\tau$ is completely independent of every other
state at the same $\tau$, so they can be solved in parallel.  After each quarter is
done, we seed the next quarter from the quarter boundary values and repeat.

Not every state in the theoretical state space is reachable from the opening kickoff.
Before solving, we run a forward BFS from the start state to enumerate only the states
that can actually occur under any sequence of plays.  This reduces the number of states
to solve from hundreds of millions to approximately $30$ million for Q4.

\begin{algorithm}[ht]
\caption{Full-game backward induction solver}
\label{alg:solve}
\begin{algorithmic}[1]
\Require Offense chart $C$, punt chart $P$, start state $s_0$
\State $\mathcal{S} \gets \textsc{ForwardBFS}(C, P, s_0)$ \Comment{Only reachable states}
\For{$q \in \{4, 3, 2, 1\}$}
  \State Apply quarter-boundary seeding (e.g.\ halftime kickoff)
  \For{$\tau \in \{1, 2, \ldots, 60\}$}
    \For{each reachable state $s$ at $(q, \tau)$ \textbf{in parallel}}
      \If{$s$ is a scrimmage state}
        \State Build $A^{(s)}$; check for saddle point
        \State $V(s) \gets$ saddle-point value or LP value of $A^{(s)}$
      \ElsIf{$s$ is a kickoff or safety-kick state}
        \State $V(s) \gets$ closed-form expected value
      \EndIf
    \EndFor
  \EndFor
\EndFor
\Return $V$
\end{algorithmic}
\end{algorithm}

\subsection{Value Storage}

Rather than a hash map over 30 million keys, values are stored in dense multi-dimensional
NumPy arrays indexed directly by state coordinates, giving O(1) lookup:
\begin{itemize}
  \item Q3/Q4 scrimmage: shape $(2,\, 61,\, 99,\, 4,\, 30,\, 71)$, approximately
        $823$~MB.  No $h$ dimension because $h=0$ in the second half.
  \item Q1/Q2 scrimmage: shape $(2,\, 61,\, 99,\, 4,\, 30,\, 71,\, 2)$, approximately
        $1.65$~GB.
  \item Kickoff and safety-kick arrays: negligible size.
\end{itemize}
Total memory: approximately $2.47$~GB.  Unsolved slots are initialised to \texttt{NaN};
any read of an unsolved slot raises an error immediately, catching dependency bugs.

%────────────────────────────────────────────────────────────────────────────────
\section{Computational Implementation}
\label{sec:implementation}
%────────────────────────────────────────────────────────────────────────────────

Solving ${\sim}30$ million matrix games is not expensive in principle --- each game is
tiny (${\leq}22\times10$) --- but doing it 16 times over (once per timeout-count
combination) requires care.  This section describes the engineering decisions that
make the solve tractable.

\subsection{Score Differential Clipping}

The score differential $\delta$ (offense minus defense) is theoretically unbounded,
but a lead of more than 35 points in a single quarter is effectively insurmountable.
We clip $\delta$ to the range $[-35, 35]$, giving 71 distinct values.  This is not
just an approximation made for convenience: the terminal utility $U(\delta) = \operatorname{sgn}(\delta)$
is flat outside $\{-1,0,+1\}$, so states with $|\delta| > 35$ already have a certain
outcome and their values propagate trivially.  Clipping makes the state space finite
and dense-array-representable without affecting the equilibrium strategies in any
reachable game situation.

\subsection{Symmetric Possession Representation}

Rather than tracking both field position from Team~1's perspective \emph{and} which
team has the ball, we always write the state from the current offense's point of view.
Field position $x$ is always yards from the \emph{current offense's} own goal.  When
possession changes, we flip $x \to 100 - x$, $\delta \to -\delta$, $h \to -h$, and
negate the continuation value.  This halves the state space and means a single value
table covers both teams --- there is no asymmetry to track.

\subsection{Pure-Strategy Saddle-Point Shortcut}

Before calling the LP solver on a payoff matrix $A$, we check whether a saddle point
exists.  A saddle point is a position in the matrix that is simultaneously the smallest
value in its row and the largest value in its column --- in other words, a play the
offense always wants to run regardless of what the defense calls, and a card the defense
always wants to call regardless of what the offense runs.

Concretely: if the maximum of the row-minimum values equals the minimum of the
column-maximum values,
\[
  \underbrace{\max_a \min_b A_{a,b}}_{\text{maximin}} = \underbrace{\min_b \max_a A_{a,b}}_{\text{minimax}},
\]
then the saddle point is a pure Nash equilibrium and no LP is needed.  The game value
is exactly the saddle-point entry.  This check costs $O(22 \times 10)$ arithmetic
operations, compared to the LP which requires hundreds of simplex iterations.  In
practice it fires on roughly 40--60\% of game states (most clearly in short-yardage
situations and end-of-game scenarios), roughly halving total solve time.

\subsection{Fork-Based Parallelism}

Two levels of parallelism are exploited.

\paragraph{Within a $\tau$ layer.}
States at the same $(q, \tau)$ are completely independent --- they read values only
from already-solved lower-$\tau$ states.  The implementation distributes them across
CPU cores using Python's \texttt{multiprocessing.Pool} with the \texttt{fork} start
method.  Forking is critical: it makes the full ${\sim}2.5$~GB value table available to
every worker at zero serialization cost (the operating system shares pages
copy-on-write), so workers pay only for the small number of pages they actually
write.

\paragraph{Across timeout-count combinations.}
Combinations $(\theta_1, \theta_2)$ with the same total $\theta_1 + \theta_2$ are
independent of each other --- they read only from already-solved combinations with
smaller totals.  Rather than solving them sequentially, we fork one process per
combination at each level and solve them simultaneously.  With 16 workers and 4
combinations at the most populated level, this gives up to $4\times$ speedup on the
timeout solve.  Workers are split evenly: if $k$ combinations run in parallel each
gets $\lfloor n_\text{workers}/k \rfloor$ state-level workers.

\subsection{Runtime and Memory}

The Q4 base solve (no timeouts remaining) runs in approximately 8 hours on a 16-core
machine.  With the saddle-point shortcut and parallel combo execution the full 16-table
solve is estimated to take 1--2 days on the same hardware.  Peak RAM usage is
approximately $3 \times 823\,\text{MB} \approx 2.5\,\text{GB}$ for the active value
tables at any one time; completed tables are saved to disk and evicted.

%────────────────────────────────────────────────────────────────────────────────
\section{Timeout Decisions}
\label{sec:timeouts}
%────────────────────────────────────────────────────────────────────────────────

Each team is entitled to three timeouts per half.  Calling a timeout after a play reduces
the clock charge according to Table~\ref{tab:to}.

\begin{table}[ht]
\centering
\caption{Clock time counted after a timeout is called, by original play duration.}
\label{tab:to}
\begin{tabular}{ccc}
\toprule
Original play duration & Ticks without TO & Ticks after TO \\
\midrule
45 sec (long gain)      & 3 & 1 \\
30 sec (normal play)    & 2 & 0 \\
15 sec (OOB / penalty)  & 1 & 0 \\
\end{tabular}
\end{table}

\noindent A timeout can save clock ticks after any play type.  The most critical case
is a 1-tick play (out-of-bounds, incomplete, or penalty) when the clock is about to
expire: calling a timeout reduces the charge to 0 ticks, preserving the remaining
time entirely and allowing one more snap.  After a 2-tick play the savings are even
larger (2 ticks saved), and after a 3-tick long gain the charge drops from 3 to 1 tick.

\subsection{The Post-Play Timeout Sub-Game}

After each play completes with tick cost $k$, either team may call a timeout to reduce
the clock charge to $k' < k$.  Let
\begin{align*}
  a &= V^{(\theta_1-1,\,\theta_2)}(s',\,k')   \quad \text{(offense calls)},\\
  b &= V^{(\theta_1,\,\theta_2-1)}(s',\,k')   \quad \text{(defense calls)},\\
  c &= V^{(\theta_1,\,\theta_2)}(s',\,k)       \quad \text{(neither calls)}.
\end{align*}
Because this is a zero-sum game, reduced clock time can benefit at most one team, so
both calling simultaneously is never rational.  Since at most one team will call, the
order of decisions is irrelevant: whichever team benefits simply calls, and the other
does not.  The sub-game value is therefore
\begin{equation}
  \label{eq:to_subgame}
  V_{\text{TO}} = \max\!\bigl(a,\;\min(b,\,c)\bigr),
\end{equation}
which replaces the plain $V^{(\theta_1,\theta_2)}(s',k)$ lookup in the outer LP.

\subsection{Layered Solve Architecture}

Adding timeouts to the state tuple would multiply the state space by $4^2 = 16$.
Instead, we maintain 16 separate value tables $V^{(\theta_1,\theta_2)}$, each with the
same array layout as the no-timeout base ($\approx 823$~MB per table, $\approx 13$~GB
total).  The tables are solved in order of increasing $\theta_1 + \theta_2$:
\[
  (0,0) \;\to\; (1,0),\,(0,1) \;\to\; (2,0),\,(1,1),\,(0,2) \;\to\; \cdots \;\to\; (3,3).
\]
The $(0,0)$ table (no timeouts remaining) is solved first and serves as the base case.
Each subsequent $(\theta_1,\theta_2)$ requires a full backward-induction pass, with the
post-play sub-game looking up already-solved predecessor tables.  The game-relevant
output is $V^{(3,3)}$: both teams begin with three timeouts.

%────────────────────────────────────────────────────────────────────────────────
\section{Comparison to McCulloch's Method}
\label{sec:comparison}
%────────────────────────────────────────────────────────────────────────────────

Table~\ref{tab:compare} summarizes the principal differences between McCulloch's approach
and ours.

\begin{table}[ht]
\centering
\caption{McCulloch's heuristic agent vs.\ the exact backward-induction solver.}
\label{tab:compare}
\setlength{\tabcolsep}{6pt}
\begin{tabular}{lll}
\toprule
Aspect & McCulloch (2004) & This work \\
\midrule
\textbf{Matrix entry value} &
  Heuristic bonus on raw yardage &
  Exact continuation value $Q_s(a,b)$ \\[2pt]
\textbf{Optimality guarantee} &
  Local Nash per play; global not guaranteed &
  Global Nash for the full game \\[2pt]
\textbf{Clock} &
  Not modeled explicitly &
  Exact 15-sec tick accounting, \\
  & & two-minute warning, timeouts \\[2pt]
\textbf{Fourth-down decisions} &
  Fixed heuristic cutoffs &
  Optimal from $V(s)$ for punt/FG/go \\[2pt]
\textbf{Score sensitivity} &
  Bonus scaling by score margin &
  Terminal utility $U(\delta)\in\{-1,0,+1\}$ \\[2pt]
\textbf{Field goal} &
  Heuristic range threshold &
  Exact probability by distance \\[2pt]
\textbf{Kickoff / safety kick} &
  Not modeled as strategic decisions &
  Full stochastic transitions; \\
  & & onside kick option included \\[2pt]
\textbf{Punts} &
  Fixed utility for receiving position &
  Optimal punt vs.\ go-for-it trade-off \\[2pt]
\textbf{Timeouts} &
  Not modeled &
  Pure-strategy sequential sub-game; \\
  & & 16-table layered solve \\[2pt]
\textbf{State space solved} &
  Each play independently &
  ${\sim}30$M states $\times$ 16 TO combos \\[2pt]
\textbf{Tunable parameters} &
  Several heuristic constants &
  None \\
\bottomrule
\end{tabular}
\end{table}

The most consequential difference is the matrix entry.  In McCulloch's formulation,
$A^{(s)}_{a,b}$ is an approximate scalar that proxies game value.  In ours, it is
$Q_s(a,b) = \mathbb{E}[V(s') \mid s,a,b]$, which is the exact expected value of the
resulting position under equilibrium play from that position onward.  Consequently, our
equilibrium mixed strategy is a Nash equilibrium for the \emph{game of Football Strategy}
itself, not merely for the local matrix game with surrogate payoffs.

\paragraph{On the correctness of heuristic strategies.}
It is in principle possible that McCulloch's heuristics happen to produce strategies close
to the true equilibrium; whether this is so is an empirical question that our solver can
now answer directly, by comparing the heuristic mixed strategies to the exact equilibrium
strategies at each state.

%────────────────────────────────────────────────────────────────────────────────
\section{Red-Zone Play Restrictions and Chart Parsing}
\label{sec:redzone}
%────────────────────────────────────────────────────────────────────────────────

The Ball Control Offense Plays Chart restricts legal play sets near the opponent's goal
line.  In offense-relative coordinates:

\begin{itemize}
  \item Plays 17--20 are illegal when $x \geq 80$ (within the opponent's 20).
  \item Plays 13--16 are illegal when $x \geq 90$ (within the opponent's 10).
\end{itemize}

The chart is parsed cell-by-cell into a typed intermediate representation before solving.
Each cell token is classified as one of: \textsc{Yards}, \textsc{Incomplete},
\textsc{Fumble}, \textsc{LongGain}, \textsc{Interception}($r$), \textsc{Penalty}($y$),
\textsc{LossOfDown}($y$), or an OR-choice \textsc{Choice}$(\text{branch}_1,
\text{branch}_2, \text{resolver})$ where \texttt{resolver} $\in \{\text{offense},
\text{defense}\}$ indicates which team selects the better branch.  This separation between
parsing and solving simplifies both the transition engine and the LP construction.

A subtle rule applies to interceptions with large negative return yards (the intercepting
defender runs toward their own end zone): if the resulting field position exceeds the
physical boundary of the defender's end zone (more than 110 yards from the offense's own
goal line in offense-relative coordinates), the play is treated as an incomplete pass
rather than a touchback.

%────────────────────────────────────────────────────────────────────────────────
\section{Scale and Computational Results}
\label{sec:results}
%────────────────────────────────────────────────────────────────────────────────

The fourth-quarter solve enumerates approximately $30.9$ million reachable states across 60
clock layers via forward BFS.  Of these, approximately 615,000 are unique
\emph{scrimmage core} states (the $(x,d,y,\delta,h)$ component), with each appearing at
multiple clock values.  Each scrimmage state requires solving a matrix LP with at most 22
rows (offense actions) and 10 columns (defense cards).

The no-timeout base ($\theta_1=\theta_2=0$) runs in approximately 8 hours on an
8-core machine; the full 16-table solve requires approximately $16\times$ that time.
Each value table is stored in approximately $823$~MB as a dense NumPy array.

The value $V(s_0)$ at the initial kickoff state $(q=4,\tau=60,\delta=0)$ gives the exact
win probability for the team kicking off at the start of the fourth quarter under a tied
score.  Any deviation from the equilibrium mixed strategies by either player can only
decrease that player's win probability.

%────────────────────────────────────────────────────────────────────────────────
\begin{thebibliography}{9}

\bibitem{mcculloch}
Sean McCulloch.
\newblock A game-theoretic intelligent agent for the board game \emph{Football
  Strategy}.
\newblock \emph{Unpublished manuscript}, 2004.

\bibitem{neumann}
John von Neumann.
\newblock Zur Theorie der Gesellschaftsspiele.
\newblock \emph{Mathematische Annalen}, 100(1):295--320, 1928.

\bibitem{shapley}
Lloyd Shapley.
\newblock Stochastic games.
\newblock \emph{Proceedings of the National Academy of Sciences}, 39(10):1095--1100, 1953.

\bibitem{puterman}
Martin Puterman.
\newblock \emph{Markov Decision Processes: Discrete Stochastic Dynamic Programming}.
\newblock John Wiley \& Sons, 1994.

\bibitem{highs}
Julian Hall, Ivet Galabova, Leona Gottwald, and Michael Feldmeier.
\newblock {HiGHS}: High performance software for linear optimization.
\newblock \url{https://highs.dev}, 2023.

\end{thebibliography}

\end{document}
